%% ============================================================================
%% sections/00-abstract.tex -- Abstract
%% Class-independent section file: it is \input by main-acm.tex and main-ieee.tex.
%% Rules (see README-build.md): no class-specific commands; cite with \cite only;
%% figure sizes come from the float budget and table bodies from tables/, never from this file.
%% Every paragraph and every caption is one numbered unit that starts with a
%% marker line "%% P:<id>".  Ids are frozen: never renumber, never reuse.
%% The same ids appear in paper.zh-CN.md; check_align.py compares the two.
%% Word budget of this section: 180 words of running text.
%% ============================================================================
%% (no blank lines in this file: it is read inside the abstract environment)
%% H1: Abstract
%% P:abs-p1
\ptag{abs-p1}%
A GEMM result can be consumed where it is produced, in a fused epilogue, or take a scalar or \blockscaled form and pass to a separate vector unit.
The choice changes where rounding happens, how long execution resources stay occupied, and how much intermediate state is stored and in what form.
We study when keeping the result \blockscaled across this matrix-to-vector \handoff pays off, in a processing element (PE) pairing a systolic \TensorUnit with two \VectorCore{}s \TBD{D1}{form, width and division of labor of the two Vector cores}.
We compare five placements of the \handoff: up to four as compile-time configurations of one RTL, the fused one with a \evSlice and a calibrated \evModel; numerical error is evaluated on synthetic inputs.
We characterize where this pays off by fusion feasibility, result lifetime and metadata transaction granularity \TBD{E4}{where the block form pays off and where it does not}, and report the cost of a \failvisible \handoff \TBD{E8}{implementation cost} and a system-level projection \TBD{P}{system-level result; source of performance data}.
Decoding with three public models from two generations of one family (DeepSeek-V4-Pro and V4-Flash, then V4.1-Flash) motivates the design and is the evaluation workload.
Today's RTL prototype has one \VectorCore and implements one placement; physical results are \evPostRoute{}s on the ASAP7 predictive PDK, not silicon measurements.
